% =====================================================================================
%  Chapter 3 — System Modelling and Residual-Force Formulation
%  Thesis: Comparison of Control Strategies for Interaction-Force Aware Multirotor Teams
%
%  DRAFT, 23 August 2026 (writing-track item W.6).
%  REVISED 2026-10-14 (thesis-revision run, Phase 4). Markers: % NEW, % CHANGED, % DECISION,
%  % VERIFY, % REMOVED. Plan rows P3-xx refer to docs/thesis_revision/stage3_revision_plan.md.
%  REVISED 2026-10-15 (round 2: final study controller; D-01, D-02, D-04, D-06, D-11 resolved by the
%  author). Code basis now: crazyflie-firmware controller_omar_indi.c, ns2_residual.c,
%  platform_defaults_cf21bl.h; tools/residual (training, offline labels).
%
%  SOURCES.
%  Claims about prior work: verify against the paper itself, never against a summary or ledger
%  in docs/ (see the ch2 header). This file is the source of truth for what it states.
%  Claims about OUR system: verifiable against the code listed above (round 2). Where this chapter states a
%  constant or an equation, it is the one in the code, not an idealisation of it. If the code
%  changes, this chapter changes.
%
%  FRAME CONVENTION: z-up (ENU), gravity acts along -e_3. This matches the implementation
%  and Crazyswarm2. Note that docs/04_Unified_Residual_Wrench_Model.md was written with a
%  +g e_3 term; that document predates this chapter and is superseded on this point.
%
%  Citation baseline: 2026-08-27 (see docs/thesis/CITATION_BASELINE.md).
%  After that date: only NEW cites or EDITED prior-work sentences need re-checking.
% =====================================================================================

\chapter{System Modelling and Residual-Force Formulation}
\label{ch:model}

Chapter~\ref{ch:related} surveyed families that measure a residual, predict one, plan over one, or
learn a corrective action with no force estimate at all. Those families are comparable as
\emph{controllers}, on the same vehicles, trajectories and formations. They are not required to
compute the same internal quantity.

% CHANGED (P3-24): the prediction is the vertical component of the same quantity.
% CHANGED (round 2): one definition for all strategies; computed onboard or offline.
This chapter defines the residual acceleration $a_{\mathrm{res}}$ that \emph{this} study measures
and records. It has one definition for every strategy: the controller computes it onboard when it
uses it, and otherwise it is computed after the flight from the recorded signals in the same way.
The strategies therefore refer to one quantity rather than to several differently-defined ones:
Strategies~1 and~3 use the measured residual in all three axes, and the learned model of
Strategies~2 and~3 predicts its vertical component. A predicted residual can therefore be checked
against a measured one on the same flight. That is how a learned model is evaluated here.

The chapter also states the assumptions the formulation rests on, and where each one is expected
to fail. Several are known to be imperfect, the summation in
Section~\ref{sec:model-deepsets} most conspicuously, and stating them here is what allows
the discussion chapter to interpret a negative result as evidence about an assumption
rather than as an unexplained failure.

% -------------------------------------------------------------------------------------
\section{Notation and conventions}
\label{sec:model-notation}

Vectors are expressed in a world frame that is \emph{east-north-up}: the third basis vector $e_3$
points away from the ground and gravity acts along $-e_3$. This matches the convention used by the
flight software and the motion-capture system, and is stated explicitly because the opposite
convention is common in the quadrotor literature and the sign of the gravity term differs between
them.

% CHANGED (P3-18): vee map, units of rotor speed, the set k_t, and the two uses of the hat.
For vehicle $i$: $p_i \in \mathbb{R}^3$ is position, $v_i \in \mathbb{R}^3$ velocity,
$R_i \in SO(3)$ the rotation from body to world, and $\omega_i \in \mathbb{R}^3$ the body angular
velocity. The hat map $\hat{\cdot} : \mathbb{R}^3 \to \mathfrak{so}(3)$ is defined by
$\hat{a}b = a \times b$, and the vee map $(\cdot)^\vee$ is its inverse, so that
$(\hat a)^\vee = a$. A hat over an angular velocity, as in $\hat\omega_i$, always denotes the hat
map; a hat over an acceleration, as in $\hat a_{\mathrm{res}}$, denotes a \emph{prediction}. The
vehicles are homogeneous, of mass $m$ and diagonal inertia $J$. Rotor $k$ spins at $\varpi_k$, in
radians per second, and produces thrust $k_t\varpi_k^2$ along the body $z$-axis, with one thrust
constant $k_t$ for all four rotors.
% CHANGED (round 2): one thrust constant (MOTORRPM2FORCE, platform_defaults_cf21bl.h:65); the sensor
%   reports revolutions per minute, converted to rad/s (controller_omar_indi.c L188-191).

% -------------------------------------------------------------------------------------
\section{Rigid-body model with a residual wrench}
\label{sec:model-dynamics}

Each vehicle is modelled as a rigid body driven by a collective thrust $f_i$ along its body
$z$-axis and a body torque $\tau_i$:
%
\begin{align}
  \dot{p}_i &= v_i, &
  m\,\dot{v}_i &= -m g e_3 + f_i R_i e_3 + f_{\mathrm{res},i}, \label{eq:model-trans}\\
  \dot{R}_i &= R_i \hat{\omega}_i, &
  J\dot{\omega}_i &= -\omega_i \times J\omega_i + \tau_i + \tau_{\mathrm{res},i}.
  \label{eq:model-rot}
\end{align}
%
The terms $f_{\mathrm{res},i}$ and $\tau_{\mathrm{res},i}$ form the \emph{residual wrench}: by
construction they collect every force and moment that \eqref{eq:model-trans}--\eqref{eq:model-rot}
otherwise omits. This includes blade flapping, rotor drag, ground effect and errors in $m$, $J$ and
$k_t$, as well as the inter-vehicle aerodynamic interaction that is the subject of this thesis.

That the residual is a catch-all rather than a physical quantity has two consequences that matter
for the experiments. First, a residual measured in isolated free flight is not zero; it is whatever
the nominal model gets wrong about a single vehicle. Second, the interaction can only be isolated
by \emph{comparison}: the residual measured in formation is compared with the residual measured
without a neighbour. The coplanar control scenarios in Chapter~\ref{ch:setup} exist for exactly
this purpose.

\paragraph{Scope: the force is predicted, the moment only measured}

% CHANGED (round 2): switch B compensates a measured moment residual (Ch4 Sec. ctrl-s1).
The learned model of this thesis predicts the \emph{force} residual $f_{\mathrm{res}}$, and only its
vertical component. The moment residual $\tau_{\mathrm{res}}$ is not predicted. The strategies that
use the measurement, Strategies~1 and~3, correct the attitude loop with a measured moment term
(Chapter~\ref{ch:control}); the others leave the moment to the attitude loop, which rejects it as a
disturbance in the usual way. This is a scoping decision, not a claim that $\tau_{\mathrm{res}}$ is
negligible.

The decision has company. Neural-Swarm compensates the vertical interaction
force~\cite{shi2020neuralswarm}, and the hardware controller of Hsieh et
al.~\cite{hsieh2026flatness} likewise commands thrust and angular rates, leaving residual torque to
the rate loop. Kiran et al.~\cite{kiran2025downwash} and Gielis et al.~\cite{gielis2023aggregate}
do measure both force and moment, but on a rig rather than as a flight residual.

We therefore do not learn $\tau_{\mathrm{res}}$. The cost is that, without the measured moment term,
interaction torque is left to the attitude loop. The benefit is a single learned output that fits
the onboard budget. The discussion chapter returns to what leaving the moment unpredicted may cost.

% -------------------------------------------------------------------------------------
\section{Measuring the residual}
\label{sec:model-measure}

The residual can be recovered without any model of the interaction, by comparing what the vehicle
actually accelerated with what the nominal model says it should have accelerated given the
actuation actually applied. Rearranging \eqref{eq:model-trans},
%
\begin{equation}
  \frac{f_{\mathrm{res},i}}{m}
  \;=\; \underbrace{\dot{v}_i}_{a_{\mathrm{meas}}}
      - \underbrace{\left( \frac{f_i}{m} R_i e_3 - g e_3 \right)}_{a_{\mathrm{model}}} .
  \label{eq:model-ares}
\end{equation}
%
% CHANGED (P3-10): the accelerometer statement made precise; implementation details from the code.
Both terms are available onboard. An accelerometer does not measure $\dot v_i$ directly: it
measures the \emph{specific force}, the acceleration of the body minus the acceleration of gravity,
expressed in the body frame, $\tilde a_i = R_i^\top(\dot v_i + g e_3)$. A vehicle resting on the
ground therefore reads $\tilde a_i = R_i^\top g e_3$, one $g$ upwards, although it does not
accelerate. Rotating the reading into the world frame and subtracting $g e_3$ gives
$a_{\mathrm{meas}} = R_i\tilde a_i - g e_3 = \dot v_i$. The model acceleration requires the applied
thrust, which is reconstructed from measured rotor speeds as $f_i = \sum_{k=1}^{4} k_t\varpi_k^2$.
The rotor speeds come from an optical rotor-speed sensor on each vehicle. Both accelerations, the
rotated accelerometer reading and the model acceleration, pass identical second-order Butterworth
low-pass filters at \SI{30}{\hertz}, sampled at \SI{500}{\hertz}, before they are subtracted.
% CHANGED (round 2; D-06 resolved): code, controller_omar_indi.c L179-196 (rotor speeds from the
%   optical deck, log variables rpm.m1..m4, raw), L250-262 (both accelerations filtered, 30 Hz,
%   L131-136). The DShot telemetry and the plausibility filter of the retired implementation are not
%   used for control.

% NEW (P3-15): a worked example with illustrative values, so the reader sees the size of a_res.
\paragraph{Two illustrative magnitudes.} The numbers in this paragraph are illustrative and are
not measurements of this project. Consider a vehicle of $m = \SI{0.041}{\kilo\gram}$ in hover, so
that the true thrust equals $mg \approx \SI{0.40}{\newton}$ and $\dot v_i = 0$. If the identified
thrust constant overstates the real thrust by \SI{5}{\percent}, the model believes the rotors
deliver $1.05\,mg$, so $a_{\mathrm{model}} = +0.05\,g\,e_3$ and \eqref{eq:model-ares} returns
$a_{\mathrm{res}} \approx -0.49$\,\si{\metre\per\second\squared}, a steady downward ``force'' that
no neighbour produced. For comparison, Kiran et al.\ measured on a static rig that the lower of two
vertically stacked Crazyflies, \SI{13}{\centi\metre} apart and driven at a constant hover voltage,
produced about \SI{65}{\percent} of its nominal hover thrust~\cite{kiran2025downwash}. If a loss of
that size were not captured by the rotor-speed thrust model, it would appear as a residual of the
order of $0.35\,g$, about \SI{3.4}{\metre\per\second\squared}. The two magnitudes differ by less
than an order of magnitude, which is why model error must be bounded separately
(Section~\ref{sec:model-measure-fail}).

% CHANGED (round 2): computed onboard with switch A, otherwise after the flight; one definition.
When the measured residual is used (switch~A, Chapter~\ref{ch:control}), the controller evaluates
\eqref{eq:model-ares} at its own rate of \SI{500}{\hertz}. When it is not used, the controller does
not compute it, and the same quantity is computed after the flight from the recorded accelerometer
signal, attitude and rotor speeds, with the same constants and the same filters. The definition is
therefore one for all strategies. This matters twice. The training data of the learned model are
flown under the \emph{uncompensated} controller, so their labels must exist when no INDI term is
active; and comparing a predicted residual to a measured one on the same flight requires both
quantities for every flight. That is how a learned model is evaluated.
% CHANGED: code, controller_omar_indi.c L231-262 (residual filters updated only when indi & 1);
%   tools/residual/omar_labels.py (offline reconstruction with the controller's filter chain);
%   tools/usd_study_config.txt (500 Hz; rpm.m1..m4, acc.x/y/z, quaternion recorded).
% VERIFY: rotor-speed source of the offline labels (optical sensor as in the controller, or motor
%   telemetry) is open (round2_open_items.md).

\subsection{What this measurement requires, and how it fails}
\label{sec:model-measure-fail}

Four properties of \eqref{eq:model-ares} shape the experimental design.

\paragraph{It requires rotor-speed sensing.} Without $\varpi_k$ the applied thrust is unknown and
$a_{\mathrm{model}}$ cannot be formed. This is the hardware prerequisite identified in
Section~\ref{sec:rw-reactive}, and the same one that Cobo-Briesewitz et
al.~\cite{cobobriesewitz2026lindi} remove by learning the residual instead of computing it.
% NEW (P3-14): evidence from the same paper on the same airframe family.
How much the quality of the thrust estimate matters is visible in their own results: on a
Crazyflie~2.1, an INDI variant that reconstructed thrust from the motor command rather than from
measured rotor speed tracked worse than the uncompensated geometric controller, which the authors
read as INDI's dependence on accurate sensor measurements~\cite{cobobriesewitz2026lindi}.
Without the rotor-speed sensor the controller does not form \eqref{eq:model-ares} at all, and the
switches that use it have no effect; the residual terms then read \emph{identically zero}. That
failure is silent and is indistinguishable, in a log file, from a flight in which no interaction
occurred. Verifying that the rotor speeds and $a_{\mathrm{res}}$ are non-zero is therefore the first
check of every experimental session (Chapter~\ref{ch:setup}).
% CHANGED: code, controller_omar_indi.c (rpm_deck_available = deck.bcRpm == 1 gates every INDI term).

% CHANGED (P3-12): delay made concrete, and the matched-filter requirement stated with the code's choice.
\paragraph{It is delayed.} $a_{\mathrm{meas}}$ comes from a noisy accelerometer signal
and must be filtered; the reconstructed thrust inherits the rotor-speed sensing delay. The estimate
therefore exists only after the disturbance has begun to act. No amount of tuning removes this: it
is the structural property that motivates prediction. Filtering costs time in a predictable way: a
second-order Butterworth low-pass filter with cut-off $f_c$ delays slowly varying signals by
$\sqrt{2}/(2\pi f_c)$, about \SI{7.5}{\milli\second} at $f_c = \SI{30}{\hertz}$ (a textbook value,
given for scale). The INDI literature adds a second requirement: the measured acceleration and the
thrust it is compared with must pass \emph{identical} filters, so that both carry the same lag and
the difference is not corrupted by a phase mismatch~\cite{tal2018indi,smeur2017cascadedindi}. Here
both pass the same second-order Butterworth filter at \SI{30}{\hertz}, designed for and run at the
controller's \SI{500}{\hertz}, so the delay figure above applies to both.
% CHANGED (round 2; D-06 resolved): code, controller_omar_indi.c L131-136 (cutoff 30 Hz, sample time
%   1/ATTITUDE_RATE) and L156 (executed at ATTITUDE_RATE = 500 Hz): nominal = realised design rate.

\paragraph{It absorbs model error.} Any error in $m$ or $k_t$ appears in $a_{\mathrm{model}}$ and is
attributed to $f_{\mathrm{res}}$. A miscalibrated thrust constant produces a residual that looks
like a constant interaction force, as the first example above shows.
% NEW (P3-13): accelerometer bias, from the cascaded-INDI literature.
An accelerometer bias enters $a_{\mathrm{meas}}$ in the same way; Smeur et al.\ note that an INDI
position loop is sensitive to such a bias, because a measurement offset becomes an offset in the
achieved acceleration~\cite{smeur2017cascadedindi}.
This is why the airframe constants are identified once and
frozen, and why the null scenarios matter: a residual that persists when no neighbour is present is
model error, not interaction.

% NEW (structure): the fourth property was implicit in the original list.
\paragraph{It is a measurement of the force, not of the controller's success.} Because
$a_{\mathrm{model}}$ is formed from the thrust the rotors actually produce, as reconstructed from
their measured speeds, a controller that compensates the disturbance changes $f_i$ and therefore
$a_{\mathrm{model}}$, but not the difference \eqref{eq:model-ares} at a given state. The same
property is stated by Shi et al.\ for their identical definition: compensating a learned force
``will only affect'' the commanded thrust ``and not the observed'' interaction
force~\cite{shi2020neuralswarm}. What compensation does change is the state the vehicle visits,
and with it the residual that state produces.
% CHANGED (round 2; D-04 resolved): Ch4/Ch5 now give the reason for dedicated training flights
%   (no evaluation on training data); the property above is consistent with it.

\subsection{Sign convention in the control law}
\label{sec:model-sign}

The compensation enters with a sign that is easy to get wrong and whose error is not obvious in
flight. From \eqref{eq:model-trans}, holding a desired acceleration $a_{\mathrm{des}}$ (the acceleration the
trajectory demands, without gravity and without the residual) requires
%
\begin{equation}
  f_i R_i e_3 \;=\; m\,a_{\mathrm{des}} + m g e_3 - f_{\mathrm{res},i},
  \label{eq:model-sign}
\end{equation}
%
% NEW (P3-16): the one substitution, and the same sign in the cited literature.
obtained by setting $\dot v_i = a_{\mathrm{des}}$ in \eqref{eq:model-trans} and solving for the
thrust vector. The desired-acceleration vector passed to the thrust-and-attitude allocation
therefore carries $-f_{\mathrm{res},i}/m$. The same sign appears, written in each paper's own frame,
in the controllers that compensate a learned or measured force: Neural-Swarm and Neural-Swarm2
subtract the predicted force from the desired
force~\cite{shi2020neuralswarm,shi2022neuralswarm2}, Tal and Karaman's linear-acceleration INDI
increment is algebraically the commanded acceleration, minus gravity, minus the estimated external
acceleration~\cite{tal2018indi}, and the geometric controller of Cobo-Briesewitz et al.\ and the
flatness-based controller of Hsieh et al.\ subtract the residual before forming the
thrust~\cite{cobobriesewitz2026lindi,hsieh2026flatness}.
The residual is \emph{subtracted}, whether it was measured
(Section~\ref{sec:model-measure}) or predicted (Section~\ref{sec:model-deepsets}), and the flown
controller does so by construction: it has no configurable sign (Chapter~\ref{ch:control}).
% CHANGED (round 2; D-01 resolved): code, controller_omar_indi.c L248 (a_d - hold) and L260, L268
%   (thrust = m (a_d + a_rpm_f - a_imu_f) . R e3, i.e. a_d - a_res_f).

% CHANGED (P3-16): the single-vehicle remark made consistent with Sec. 3.2 (the isolated residual is not zero).
Adding it instead does not merely fail to compensate: it doubles the disturbance, because the
controller then commands a thrust error equal and opposite to the correction it should have made.
The resulting behaviour is a stable, well-controlled vehicle that tracks slightly worse in
formation. An uncompensated vehicle can look the same. In single-vehicle flight the error is
hard to see, since without a neighbour the residual consists only of the vehicle's own small model
error, whose effect is doubled as well but remains small. With the sign
correct, position-INDI compensation must \emph{reduce} the formation-induced displacement. With
the sign inverted, the displacement would be larger, approximately twice the uncompensated value
in a quasi-static stack in the linear regime. With thrust saturation or filter delay the factor
need not be exactly two.
% CHANGED (round 2): "The two-robot results chapter can check that prediction" removed: the flown
%   controller has no sign parameter, so no inverted-sign flight exists.

% NEW (P3-16): the "twice" made explicit, with the assumption it rests on.
The factor of two follows from one line of arithmetic. Take the vertical axis alone, a
proportional-derivative position loop with stiffness $k_p$ and no integral action, and a constant
residual acceleration $d$ pressing the vehicle down. At rest the loop balances the residual with a
position error: $k_p e = d$ without compensation, so $e = d/k_p$; $k_p e = d - d = 0$ with the
measured residual subtracted; and $k_p e = d + d$ with it added, so $e = 2d/k_p$. For example, with
$d = 1$\,\si{\metre\per\second\squared} and $k_p = 30$\,\si{\per\second\squared} (illustrative
values) the sag is \SI{3.3}{\centi\metre}, zero, or \SI{6.7}{\centi\metre}. The argument assumes
no integral action on that axis: an integrator drives the steady-state error to zero whatever the
sign, and the difference between the signs then appears only in the transient. The flown controller
has such an integrator on the vertical axis in every strategy (Chapter~\ref{ch:control}), so for it
the argument describes the transient, not the steady state.
% CHANGED (round 2; D-02 resolved): code, vertical integral Kpos_Iz 1.5 in S0-S3.

% -------------------------------------------------------------------------------------
\section{Predicting the residual from relative state}
\label{sec:model-deepsets}

% CHANGED (P3-03): the predictor also sees the vehicle's own height and velocity.
The interaction force on vehicle $i$ is a function of the configuration of the team relative to it,
and, near the ground, of the vehicle's own height and velocity,
%
\begin{equation}
  f_{\mathrm{res},i} \;\approx\; \mathcal{F}\!\left( \{\, p_j - p_i,\; v_j - v_i \,\}_{j \neq i},\; z_i,\; v_i \right),
  \label{eq:model-F}
\end{equation}
%
where $z_i = e_3^\top p_i$ is the height above the floor. Equation~\eqref{eq:model-F} is what a
predictive strategy learns. Two structural properties are required
of any parameterisation of $\mathcal{F}$, and both follow from the physics rather than from
convenience. The value must not depend on the order in which neighbours are enumerated, and it must
be defined for a varying number of them, since a vehicle in a three-robot team has two neighbours
and one in a pair has one.

% NEW (P3-09): where the inputs come from.
The neighbours' positions reach each vehicle over the radio, with time stamps, as part of the
motion-capture broadcast that the flight stack already uses; the onboard model therefore needs no
additional communication. That stream carries positions only, so each vehicle estimates a
neighbour's velocity by finite differences of successive positions: the first sample of a neighbour
gives zero velocity, a repeated time stamp reuses the last estimate, and a gap of
\SI{500}{\milli\second} or more restarts the estimate. At most three neighbours are considered.
% CHANGED (round 2): code, ns2_residual.c:ns2_eval (peer_get_all; pt == 0 -> zero; ts_new == pt ->
%   reuse; 0 < dt < 500 ms -> difference; otherwise resync), NS2_MAX_NEIGHBOURS 3 (ns2_residual.h).

\subsection{Deep-sets parameterisation}

% CHANGED (P3-02): the equation now states what the code computes.
The parameterisation used here satisfies both properties by construction:
%
\begin{equation}
  \hat{a}_{\mathrm{res},i}
  \;=\; e_3\,\operatorname{sat}_{\bar a}\!\left( \frac{c}{m}\,
     \rho\!\left( \phi_G\!\left(-z_i,\, -v_i\right)
       + \sum_{j \in \mathcal{N}_i} \phi_S\!\left( p_j - p_i,\; v_j - v_i \right) \right) \right),
  \label{eq:model-deepsets}
\end{equation}
%
where $\phi_S$ is applied to each admitted neighbour independently, $\phi_G$ encodes the vehicle's
own height and velocity, and $\rho$ maps the summed embedding to a scalar. The set
$\mathcal{N}_i$ contains the neighbours that pass the input gate of
Section~\ref{sec:model-guards}. The scalar is the vertical component of the interaction force in
grams-force, the unit in which Neural-Swarm2 reports the learned force~\cite{shi2022neuralswarm2}; the factor $c = g/1000$ converts it to
newtons and the division by $m$ to an acceleration, and $\operatorname{sat}_{\bar a}$ limits its
magnitude to $\bar a$ (Section~\ref{sec:model-guards}). The horizontal components are zero by
construction. Summation is permutation-invariant and accepts any number of terms, so one parameter
set serves teams of different sizes. The output is an acceleration rather than a force so that it
is directly comparable with the vertical component of the measured $a_{\mathrm{res}}$ of
\eqref{eq:model-ares}; the two quantities are in the same units, in the same frame, at the same
point in the control loop.
% CHANGED: code, ns2_residual.c:ns2_eval_rel (ground term, gate, GRAMS_TO_NEWTONS, OUT_CLAMP 8).
% Source: NS2 PDF Sec. VII-C (f_a,z errors "in grams"), Fig. 5 axis |f_a| [g].

% NEW (P3-19): a toy example of the two properties and of "more general than summing forces".
\paragraph{A toy example.} The functions in this paragraph are invented for illustration and are
not the trained networks. Let a scalar embedding $\phi(\Delta z) = 1/\Delta z^2$ describe a
neighbour directly above at height difference $\Delta z$ (in metres, units dropped), and let
$\rho(s) = s/(1 + s/10)$. Two neighbours at $\Delta z = 0.3$ and $0.5$ give
$\phi = 11.1$ and $4.0$, whose sum $15.1$ is the same in either order; this is the permutation
invariance. The aggregate is $\rho(15.1) = 6.0$, whereas adding the two single-neighbour effects
would give $\rho(11.1) + \rho(4.0) = 5.3 + 2.9 = 8.2$. Because $\rho$ is nonlinear, the summed
embedding can express an aggregate that is not the sum of the pairwise effects; it cannot express
an aggregate that depends on how the two neighbours are arranged relative to \emph{each other}
beyond what their separate embeddings encode.

% CHANGED (P3-01): what is borrowed and what is ours, corrected to the code.
\paragraph{What is borrowed and what is specific to this thesis.} The \emph{family} is the permutation-invariant deep
set used by Neural-Swarm and Neural-Swarm2~\cite{shi2020neuralswarm,shi2022neuralswarm2}, and the
architecture is Neural-Swarm2's own: Neural-Swarm2 uses one $\phi$ per neighbour type, one $\rho$
per own type and treats the environment as a further neighbour, to capture ground
effect~\cite{shi2022neuralswarm2}. The weight layout keeps Neural-Swarm2's five weight sets, of
which only the small-vehicle pair and the ground branch are evaluated, because every vehicle in this
study is a Crazyflie. The layers are:
$\phi_S$ maps $6 \rightarrow 25 \rightarrow 40 \rightarrow 40 \rightarrow 20$ (relative position and
velocity of one neighbour), $\phi_G$ maps $4 \rightarrow 25 \rightarrow 40 \rightarrow 40
\rightarrow 20$, and $\rho$ maps $20 \rightarrow 40 \rightarrow 40 \rightarrow 40 \rightarrow 1$,
with rectified-linear activations in the hidden layers and a linear output, because an interaction
force may point either way. These are the widths Neural-Swarm2 reports, with a hidden dimension of
$20$~\cite{shi2022neuralswarm2}. The stored weight array has \num{19297} entries, of which \num{11461}
lie on the evaluated path; the weights are fixed in the flight software. What is specific to this
thesis is the training: the network is trained on dedicated two-vehicle training flights under Strategy~0
(Chapter~\ref{ch:control}), with a robust (Huber) loss on the
vertical component of the measured residual, and spectral normalisation, which Neural-Swarm and
Neural-Swarm2 apply~\cite{shi2020neuralswarm,shi2022neuralswarm2}, is not used. This is not the
network of Cobo-Briesewitz et al.~\cite{cobobriesewitz2026lindi}. An architecture ablation would be
a separate study and is out of scope unless time remains after the comparison.
% CHANGED: code, residual_nn.rs header and constants (L12-30, L91-125: 19297 weights; phi_S 3675,
%   phi_G 3625, rho_S 4161); Cargo.toml feature residual_nn_flash; lib.rs:rnn_service (L922-937);
%   tools/residual/train.py (target a_res_z, SmoothL1Loss, L17, L118); tools/residual/model.py
%   (no spectral normalisation). Reference code: crazyswarm2/crazyflie_sim/.../backend/neuralswarm.py.
% REMOVED (P3-01): "The layer widths and the resulting 987 parameters are ours, sized for the
%   control period of the onboard STM32 microcontroller. ... it is not a copy of Neural-Swarm2's
%   layers." and "The concrete network is small: phi maps 6 -> 16 -> 16 -> 8 and rho maps
%   8 -> 16 -> 16 -> 3 ... This is 987 parameters, evaluated onboard within the control period."
%   (describes a network the code no longer contains).

% CHANGED (round 2): evaluation rate of the final controller.
The network is not evaluated at every control step. It runs on every tenth tick of the
\SI{1}{\kilo\hertz} main loop, that is at \SI{100}{\hertz} or every fifth step of the
\SI{500}{\hertz} controller, and its output is held in between; the prediction is computed and
recorded whether or not it is used by the controller (Chapter~\ref{ch:control}). The sizing is discussed in
Chapter~\ref{ch:control}; what matters here is that \eqref{eq:model-deepsets} is the model, and the
layer widths are an implementation choice within it.
% CHANGED: code, ns2_residual.c:ns2_update (tick % rnn.div, div 10); controller_omar_indi.c L156.

\paragraph{These equations are the flown ones.} Equation~\eqref{eq:model-deepsets}, the layer widths
above, the input gate of Section~\ref{sec:model-guards} and the output bound are those of the
flight software, not an idealisation of them. The same applies to \eqref{eq:model-ares} and the
sign convention of Section~\ref{sec:model-sign}. If the flight software changes, this chapter
changes with it; a modelling chapter that has drifted from the flight software describes a system
that was never flown.
% CHANGED (round 2): file names removed from the visible text; code basis in the header comment.

\subsection{The superposition assumption}
\label{sec:model-superposition}

In words: each neighbour is turned into a short list of numbers by the network $\phi_S$, the lists of
all admitted neighbours are added to the list for the ground, and a second network $\rho$ turns the sum
into one force. The short list is called the latent space (or embedding). Equation~\eqref{eq:model-deepsets}
assumes the aggregate interaction is a sum of per-neighbour contributions in the latent space of
$\phi_S$. This is more general than summing forces directly,
$\rho$ is free to be nonlinear in the sum, but it is still an assumption, and the literature
states plainly that it does not hold exactly. Shi et al.~\cite{shi2020neuralswarm} report that with
more than two vehicles the aerodynamic effect is not a simple superposition of each pair, and Gielis
et al.~\cite{gielis2023aggregate} motivate a learned aggregation model, measured on a load stand,
precisely on the grounds that a single-vehicle downwash model is unlikely to suffice for predicting
aggregate forces.
% NEW (P3-20): two further readings from cited papers that bear directly on RQ3.
Two further observations bear directly on the transfer question of this thesis. Shi et al.\ found
that a network trained on two-vehicle data alone flew worse than the uncompensated baseline on
three- and four-vehicle swaps, ``since the training data does not include data on
superpositions''~\cite{shi2020neuralswarm}. And Hsieh et al.\ nevertheless model the aggregate on
three vehicles as a pairwise summation of a model trained on two, citing the same load-stand
measurements~\cite{hsieh2025onlineadaptation}; the linear regime that Gielis et al.\ identify is
consistent with that choice, and their nonlinear regime is not~\cite{gielis2023aggregate}.

The assumption is retained here for two reasons. It is what makes one trained model applicable to
teams of two and three without retraining, which the comparison requires. And a controller does not
need an exact model: it needs one whose residual after compensation is smaller than the residual
before it. The relevant question is therefore not whether \eqref{eq:model-deepsets} is exact. It is not.
The question is how large the resulting error is on a deployed model. The multi-robot results
chapter measures that.

\subsection{Input conditioning}
\label{sec:model-guards}

% CHANGED (P3-04): the gate of the code replaces the distance cut-off and the near-field clamp.
One gate is applied to the inputs of $\phi_S$, and it is part of the model rather than defensive
programming, because it changes the function being learned. A neighbour $j$ enters the sum in
\eqref{eq:model-deepsets} only if
%
\begin{equation*}
  \lvert x_j - x_i \rvert < \SI{0.2}{\metre}, \qquad
  \lvert y_j - y_i \rvert < \SI{0.2}{\metre}, \qquad
  \lvert v_{x,j} - v_{x,i} \rvert < \SI{1.5}{\metre\per\second}.
\end{equation*}
%
Neural-Swarm2 defines a vehicle's neighbours through an interaction volume that it chooses as ``a
cuboid based on observed interactions''~\cite{shi2022neuralswarm2}; the numerical bounds above are
those of the authors' public code, including the asymmetry that the relative
velocity is gated along $x$ only. Three consequences follow, and each matters for the experiments.
A neighbour that is offset laterally by more than \SI{0.2}{\metre} contributes nothing, at any
vertical separation: the model treats the wake as a column about \SI{0.4}{\metre} wide, which is
consistent in kind with the load-stand observation that a multirotor's downwash stays concentrated
within about a body radius over a long vertical distance~\cite{gielis2023aggregate}.
There is no vertical cut-off: a neighbour well above is admitted, and its effect is left to the
learned $\phi_S$. And there is no near-field clamp: the model is evaluated at whatever separation
occurs, so it is the scenario library and its safety limits (Chapter~\ref{ch:setup}) that keep the
inputs within the range the training data covered. A learned function is least trustworthy exactly
where the physics is strongest, and that range limit is what bounds the extrapolation.
% CHANGED: code, ns2_residual.c GATE_DXY/GATE_DVX and ns2_eval_rel; earlier: residual_nn.rs GATE_DXY/GATE_DVX (L139-140) and ResidualNet::eval (L281-286);
%   residual_nn.rs header L36-44 ("NOT a distance cutoff"; near-field rescaling removed).
%   Reference: crazyswarm2/crazyflie_sim/crazyflie_sim/backend/neuralswarm.py:78.
% REMOVED (P3-04): "A neighbour farther than 2 m is excluded from the sum. ... the null scenarios
%   C1--C3 check this." and "A separation closer than 4 cm is clamped to that value, preserving
%   direction. ..." (mechanisms the code no longer contains).

% CHANGED (P3-05): the bound applies to the scalar output.
The output is bounded as well: the magnitude of the vertical prediction is limited to
$\bar a$ = \SI{8}{\metre\per\second\squared}, comfortably above the
interaction accelerations expected between vehicles of this size, and far below what an
ill-conditioned network can emit; a non-finite output is replaced by zero and flagged in the same
way. This is not a hypothetical failure mode. In simulation, Gu et
al.~\cite{gu2025comparative} report a radial-basis-function estimator crashing the vehicle under
periodic wind, and an extreme learning machine becoming unstable on landing under ground effect;
their echo state network and Gaussian process did not. The authors say the failure \emph{may be
attributed} to the random initialisation of the input-to-hidden weights, which \emph{``can produce
excessively large prediction values''}. Their vehicle is a single quadrotor and their disturbances
are wind and ground effect, so this is not evidence about inter-vehicle interaction, but it is
evidence that an unbounded learned term can destabilise a controller. A saturating output converts
that failure into a bounded and, because saturation is logged, a visible one.
% CHANGED: code, ns2_residual.c:ns2_eval_rel (isfinite, OUT_CLAMP).

The training data must pass through the same gate as the deployed model, and the training pipeline
applies it. A model fitted to ungated inputs is fitted to a different distribution than the one it
will meet, and the discrepancy is invisible until predicted and measured residuals are compared in
flight.
% CHANGED: code, tools/residual/dataset.py:normalisation_rel (statistics over gated neighbours).

\subsection{Normalisation}

% CHANGED (P3-06): two first layers, separate statistics.
Input scaling is folded into the first layers of the network at export rather than applied
onboard. There are two such layers, the first layer of $\phi_S$ and the first layer of $\phi_G$,
and each receives its own statistics, because a neighbour's relative state and a vehicle's own
height and velocity follow different distributions. For
weights $W_1$ and bias $b_1$ trained on inputs normalised by mean $\mu$ and standard deviation
$\sigma$,
%
\begin{equation}
  W_1 \frac{x - \mu}{\sigma} + b_1
  \;=\; \underbrace{\left(W_1 \oslash \sigma\right)}_{W_1'} x
      + \underbrace{\left(b_1 - \left(W_1 \oslash \sigma\right)\mu\right)}_{b_1'},
  \label{eq:model-fold}
\end{equation}
%
with $\oslash$ denoting column-wise division. The two forms are arithmetically identical, so the
deployed network computes the same function as the trained one while the flight software carries no
normalisation constants that could fall out of step with the model.
% CHANGED: code, tools/residual/model.py:fold_normalisation (L224-241).
% NEW (P3-19): one line of intuition.
For example, a first-layer weight of $2$ applied to a height normalised with $\mu = \SI{1}{\metre}$
and $\sigma = \SI{0.5}{\metre}$ becomes the weight $4$ applied to the raw height, with $-4$ added to
the bias; the neuron receives the same input in both forms.

% -------------------------------------------------------------------------------------
\section{Assumptions and where they fail}
\label{sec:model-assumptions}

Table~\ref{tab:model-assumptions} collects the assumptions this formulation makes. Each is stated
with the condition under which it is expected to break, so that an experimental result can be
attributed to a specific assumption rather than to the model as a whole.

% CHANGED (P3-21): three rows added (finite-difference velocity, lateral gate, accelerometer bias).
\begin{table}[t]
  \centering
  \caption{Assumptions of the residual formulation, the conditions under which each is expected to
    fail, and the consequence for the measured residual $a_{\mathrm{res}}$ and the prediction
    $\hat a_{\mathrm{res}}$.}
  \label{tab:model-assumptions}
  \small
  \begin{tabular}{@{}p{0.30\linewidth}p{0.30\linewidth}p{0.32\linewidth}@{}}
    \toprule
    Assumption & Where it fails & Consequence \\
    \midrule
    Rigid body; no aeroelastic deformation
      & Not expected to fail at this scale
      & --- \\
    $m$, $J$, $k_t$ known and constant
      & Battery discharge; rotor wear; miscalibration
      & Model error appears as a constant residual; bounded by the null scenarios (second phase), provisionally by the repeat-to-repeat variation \\
% DECISION D-14 (new): provisional floor from repeat-to-repeat variation; author to confirm.
% NEW (P3, Phase 1 R14 Sec. 4.6): accelerometer bias row.
    Accelerometer unbiased
      & Sensor bias or drift
      & Appears as a constant residual; bounded by the null scenarios (second phase), provisionally by the repeat-to-repeat variation \\
    Residual is dominated by inter-vehicle interaction
      & Near the ground; near walls
      & Flight volume is bounded away from both (Chapter~\ref{ch:setup}) \\
    Neighbour relative states are known accurately
% CHANGED: code, ns2_residual.c:ns2_eval (finite-difference peer velocity).
      & Motion-capture dropout; latency; velocity differenced from broadcast positions
      & Prediction degrades; measurement is unaffected, which is itself a comparison result \\
% NEW: code, ns2_residual.c GATE_DXY, GATE_DVX (input gate).
    The interaction is negligible outside the input gate
      & Neighbours offset laterally by more than \SI{0.2}{\metre} that still disturb the vehicle
      & Prediction is zero where a residual is measured; visible as a prediction--measurement gap \\
    Aggregate interaction superposes in latent space
      & Three or more vehicles, overlapping wakes
      & \textbf{Known to be imperfect}~\cite{shi2020neuralswarm,gielis2023aggregate}; magnitude measured in the second, three-vehicle phase (multi-robot results chapter) \\
    Interaction lies within the training distribution
      & Unseen separations, faster relative motion
      & Prediction degrades while measurement does not; this is the crossover of \textbf{RQ2} \\
    Rotor-speed measurement available
% CHANGED (round 2): code, controller_omar_indi.c (optical deck, rpm_deck_available).
      & Rotor-speed sensor absent, or a rotor reading missing
      & Residual terms read zero or are wrong --- silent, and checked every session \\
    \bottomrule
  \end{tabular}
\end{table}

Two entries deserve emphasis because they are not merely caveats but hypotheses the experiments
test. The superposition assumption is known to be imperfect, and the experimental question is the
size of the penalty rather than its existence. And the training-distribution assumption is exactly
what separates the predictive family from the reactive one: a measurement has no training
distribution to leave. Cobo-Briesewitz et al.~\cite{cobobriesewitz2026lindi} observe precisely this
effect for their purely learned method in the no-payload case, attributing a drop in performance on
the circular trajectory to its being out of distribution relative to the training data.

% -------------------------------------------------------------------------------------
\section{Summary}
\label{sec:model-summary}

% CHANGED (P3-23): the prediction is the vertical component; both enter at one point with one sign.
% CHANGED (round 2).
The residual of \eqref{eq:model-trans}--\eqref{eq:model-rot} is what this study measures and
records, with one definition for every strategy: Strategies~1 and~3 use the measurement,
Strategies~2 and~3 a prediction of its vertical component, and Strategy~0 neither. The residual can be recovered online from
\eqref{eq:model-ares}, which needs no model of the interaction but requires rotor-speed sensing and
is available only after the disturbance acts; or its vertical component can be predicted from
\eqref{eq:model-deepsets}, which acts ahead of the disturbance but is valid only within its training
distribution, sees only neighbours inside its input gate, and assumes an aggregation that is known
to be approximate. Both are expressed as accelerations in the world frame, and both enter the
control law at the same point with the same sign, \eqref{eq:model-sign}.
% CHANGED (round 2; D-01 resolved): both are subtracted by construction.

That common definition, together with the common controller of Chapter~\ref{ch:control}, is what
makes the comparison in the two-robot results chapter a comparison rather than a collection of
separately-tuned controllers. Chapter~\ref{ch:control}
describes the strategies built on it, and the point at which each injects its knowledge of
the interaction.
